%%%%%%%%%%%%%%%%%%%%%%%%%%%%%%%%%%%%%%%%%%%%%%%%%%%%%%%%%%%%%%%%%%%%%%%%%%%%%%%%
%%
\cleardoubleoddpage%  Make sure to start each chapter on a new odd page
\chapter{Experiments}
\label{ch:experiments}


\section{Evaluation Protocol}
\label{sec:experiments:protocol}

% TODO (crucial): before introducing the other metrics, explicitly explain
% why per-pixel MSE is a deceptive metric for slow-moving GA pathology.
% Then: per-timestep MSE and RMSE in normalized space; Dice and IoU on the
% GA mask channel at a clinically meaningful 360-day horizon. Persistence
% baseline (repeat last observed state) as the floor any learned model must
% beat.


\section{Baselines}
\label{sec:experiments:baselines}

% TODO:
%   - Persistence (last-state repetition)
%   - Single-branch GNN (MP-PDE style, no moving mesh)
%   - Multi-branch GNN, larger parameter budget (isolate moving-mesh
%     contribution from capacity)
%   - Optional: simple CNN baseline (U-Net or similar) for context against
%     standard medical imaging approaches.


\section{Main Results}
\label{sec:experiments:main-results}

% TODO: table comparing all methods on MSE, RMSE, Dice@360d, IoU@360d.
% Figure with example rollouts: ground-truth sequence vs each baseline's
% predictions over several visits for a representative eye.


\section{Ablations}
\label{sec:experiments:ablations}

% TODO: each component on/off:
%   - DMM moving mesh vs uniform mesh
%   - correction branch vs full second branch
%   - ItpNet pre-training vs no pre-training
%   - pushforward training vs no pushforward
%   - covariates (Age, Sex) vs no covariates
%   - layer encoder $d_{\text{embed}} \in \{0, 32, 64, 128, 256\}$
%   - mask-channel loss-weight sweep.


\section{Moving Mesh Quality}
\label{sec:experiments:mesh-quality}

% TODO: std/range of cell volumes under the monitor function
% (following the MM-PDE paper's metric) to verify DMM generates physically
% sensible meshes on GA data. Visualize example moved meshes overlaid on
% sample states.


\section{Qualitative Analysis}
\label{sec:experiments:qualitative}

% TODO: per-eye rollout visualizations. Success cases (accurate lesion
% growth prediction). Failure modes (atypical progression patterns, very
% sparse visit schedules, very short vs long $\Delta t$).


\section{Computational Cost}
\label{sec:experiments:cost}

% TODO: training time, inference time per rollout step, memory footprint,
% comparison with baselines.


%%
%%%%%%%%%%%%%%%%%%%%%%%%%%%%%%%%%%%%%%%%%%%%%%%%%%%%%%%%%%%%%%%%%%%%%%%%%%%%%%%%
